\section{Scaling out and connecting to data in detail}
\label{app:scale}

This appendix gives the design, the measurement conditions and the evidence behind each
capability of \cref{sec:scale}, with the tests and result files that reproduce every number. The measurements were made on the
laptop of \cref{sec:performance}, with the machine shared and without the fixed frequency of
\cref{sec:playbooks}.

\subsection{Per-agent random streams and split runs}
\label{app:split}

\paragraph{Streams.} A model class that sets \code{agent\_streams = True} gives every agent
generators of its own: \code{agent.random} and \code{agent.rng} are seeded from the
scenario's seed sequence and the agent's \code{stream\_key()} (its \code{unique\_id} unless
overridden), in a spawn-key domain that \code{SeedSequence.spawn} never reaches. An agent
then draws the same numbers whatever the activation order and the population, and a
checkpoint continues every stream. The flag is off by default, and a model without it draws
from the model's two generators (\code{tests/core/test\_streams.py}).

\paragraph{Partitions and the contract.} \code{modelly.distributed} cuts a grid into balanced
blocks with their halos (torus-aware), a network by breadth-first banding with greedy boundary
refinement (not METIS \citep{karypis1998fast}: on a $30 \times 30$ lattice cut in four it cuts 134
edges against the 60 of the best split, 2.2 times it, and a tenth of a random assignment's
1,315, \code{paper/data/graph\_cut.json}; the test asserts at most three times the best), and
a continuous space by recursive bisection at the population's quantile. A model declares \code{partitionable = Partitionable(space, phases, reporters)}:
each phase names how far it reads (\code{radius}) and whether it moves agents. Before a phase
that reads across a border the workers exchange ghost copies of their border agents; after a
phase that moves agents, the movers migrate to their new partition with their streams, in a
deterministic order. The build refuses, with the reason, a model that draws from the model's
generators inside a phase, holds \code{Column} or dataset agents, schedules events of its own,
moves into cells with a capacity, keeps references to other agents in attributes, or creates
agents keyed on \code{unique\_id} during a phase. Backends are \code{local}, \code{mpi}
(\code{mpi4py}), \code{dask} and \code{ray}; \code{Sum}, \code{Min} and \code{Max} reporters
reduce across partitions, and a checkpoint writes one file per partition with an index of
digests that \code{Cluster.restore} continues draw for draw.

\paragraph{Measured.} On three reference models (\code{modelly/examples/distributed}: a grid
epidemic on blocks, a network SIR on graph parts, a continuous SIR on KD-tree boxes), every
reporter row and every agent's attributes, location and stream state are byte-identical to a
plain single-process run at 1, 2, 4 and 16 partitions on two seeds, with the compiled and the
pure-Python core (\code{tests/distributed/test\_independence.py}, also in the determinism
suite), under \code{mpirun} with 2 and 4 ranks on the grid and network models
(\code{test\_mpi.py}), and on two Dask workers. Weak scaling of the grid epidemic
(\code{benchmarks/distributed/weak\_scaling.py}, raw numbers in
\code{paper/data/weak\_*.json}) is in \cref{tab:scaling}: next to each split run of $K$
workers the script times $K$ independent copies of one block with no communication, which
separates what the machine gives from what the protocol keeps of it.

\begin{table}[htbp]
  \centering
  \small
  \caption{Weak scaling of the split grid epidemic under MPI, one run per worker count, 10,000
  people per worker, 20 steps, Open MPI 5.0.11 and mpi4py 4.1.2, wall-clock seconds per step after a
  warm-up step, build excluded, on the 16-core laptop of \cref{sec:performance} shared with
  other jobs. ``Load'' is the one-minute load average at the start of each run; ``machine'' is
  the efficiency of $K$ copies with no communication, timed by the slowest copy; ``protocol'' is the split run's time
  relative to the copies'.}
  \label{tab:scaling}
  \begin{tabular}{rrrrrrrr}
    \toprule
    Workers & People & Split, s/step & Copies, s/step & Efficiency & Machine & Protocol & Load \\
    \midrule
    1  & 10,000  & 0.0150 & 0.0110 & 100\% & 100\% & 74\%  & 3.4 \\
    2  & 20,000  & 0.0176 & 0.0190 & 85\%  & 58\%  & 108\% & 3.5 \\
    4  & 40,000  & 0.0491 & 0.0206 & 30\%  & 53\%  & 42\%  & 3.9 \\
    8  & 80,000  & 0.0775 & 0.0435 & 19\%  & 25\%  & 56\%  & 4.3 \\
    12 & 120,000 & 0.1389 & 0.0592 & 11\%  & 19\%  & 43\%  & 5.3 \\
    16 & 160,000 & 0.2732 & 0.1174 & 5\%   & 9\%   & 43\%  & 8.7 \\
    \bottomrule
  \end{tabular}
\end{table}

With 10,000 people per worker the split run takes 0.9 to 2.4 times as long per step as its
slowest copy, so the protocol keeps 42 to 56\% of what the machine gives from four workers on
(the 108\% at two workers, a split run faster than its copies, is within a single run's
noise). Under MPI the efficiency falls from 85\% at two workers to 30\% at four and 5\% at
sixteen, with the machine itself, the copies, at 53\% and 9\%. With 40,000 people per worker
(10 steps, loads of 11 to 16), the efficiency is 91, 68 and 45\% at 2, 4 and 8 workers against
104, 81 and 50\% for the copies, so the protocol keeps 74 to 80\%: a larger block has more
work per exchanged border. The Dask backend, which routes every exchange through its
scheduler (eight per step), measured 69, 35 and 22\% at 2, 4 and 8 workers (loads of 9 to 11),
each split step 2.1 to 4.0 times its copies: the scheduler's round trips, not the model, take
its time. Every efficiency here is a single run per worker count, not a median of repeats,
and a single run varies by tens of percent on this shared machine (the split run faster than
its copies at two workers shows it), so the table supports the shape of the curve and the
byte-identity claims, not the last digit of an efficiency.

\paragraph{Not met.} Efficiency at 70\% or more from 1 to 64 workers, the level a
cluster code would aim for, cannot be measured on one laptop, and the laptop's hybrid
processor is the main limit: eight copies of the same work with no communication at all
already run at 25\% of one (50\% with 40,000 people per worker), because beyond the four
performance cores the copies land on efficiency cores (at eight workers their step times
split into two groups, 0.022~s and 0.029 to 0.044~s). The Ray backend is written but untested (Ray does not install on CPython 3.15); MPI on
macOS and Windows is untested (\cref{sec:threats}); and the split-run suite runs on the three
reference models, not on the zoo, whose models would change their numbers if moved to
per-agent streams.

\subsection{A GPU backend for column phases}
\label{app:gpu}

\code{modelly.gpu} runs column phases (\code{DeviceTable}), cohort tables
(\code{DeviceCohortTable}, binomial thinning included) and the cell-space operators
(\code{DeviceCellIndex}: counts, neighbour sums, random-neighbour moves) on \code{numpy},
CuPy \citep{okuta2017cupy} or JAX \citep{bradbury2018jax}, chosen by \code{get\_backend} or
\code{MODELLY\_BACKEND}; a backend that does not load falls back to \code{numpy} with a
warning, and \cmd{import modelly} loads neither library. Random numbers come from
\code{PhiloxStream}, a stateless Philox4x32-10 \citep{salmon2011parallel} keyed by seed and
stream, checked against the Random123 known-answer vectors and identical on every backend;
binomial draws use inversion below a mean of 10 and Hörmann's BTRS rejection sampler above it
\citep{hormann1993generation}, both with the exact law (chi-square tests), inversion identical
across backends and BTRS equal in every test and run (it uses \code{log}, so a last-bit
difference is possible in principle). Without 64-bit types JAX draws single-precision
uniforms, so its runs differ from \code{numpy}'s in value, not in law. Populations larger
than the device run in batches, bit for bit the same as one batch.

\paragraph{Measured.} The script is \code{gpu\_speedup.py} in \code{benchmarks/playbooks/},
and the raw numbers are in \code{paper/data/gpu\_speedup.json}. The grid epidemic at 10
million agents on a $1000 \times 1000$ torus, 60 steps, seed 1, in fresh interpreters, wall
clock around the steps with the device synchronised, the build timed apart, one untimed run
first to fill CuPy's kernel cache: the CPU column variant (CPython 3.15) takes 20.72, 27.03
and 20.69~s for the steps and the CuPy variant (CPython 3.14, the GPU environment's interpreter;
the two interpreters differ, but the steps are dominated by numpy and CuPy kernels) 0.589, 0.673 and 0.622~s,
per-round ratios of 35.2, 40.2 and 33.3 and a median of 35.2 on an NVIDIA RTX~5050 laptop GPU
(8~GB, CUDA runtime 13.2 from the \code{cupy-cuda13x} 14.2.0 wheel). Run on \code{numpy} and on CuPy with the same seed, the variant
gives the same infected series for all 60 steps and the same SHA-256 of the final state, days
and row columns of all 10 million agents; on \code{numpy} its steps take 56.1~s, more than the
column variant's because Philox draws cost more than \code{numpy}'s generator, and 78 times its
0.72~s on CuPy. Building 10 million agent objects takes 25 to 33~s on either side, so build
plus steps is only 1.5 to 1.9 times faster: a GPU an order of magnitude faster than the
column form at ten million agents is what the steps show, not what a whole run shows.

\paragraph{Where it does not pay.} At the playbook's own 20,000 agents the steps take 0.14
to 0.16~s on CuPy against 0.08 to 0.14~s on the CPU (three seeds), and the insect cohorts, a
table of under 200 rows, take 1.52 to 1.53~s on the GPU for 120 days against 7 to 8~ms on the
CPU: kernel launches cost more than the arithmetic. JAX recompiles whenever a cohort table changes
length. Apple Metal was not measured (no Apple GPU was available), and against the original
CPU column variant, which draws from \code{model.rng}, the GPU variant agrees in distribution
only, as every playbook variant does.

\subsection{Differentiable models}
\label{app:diff}

\code{modelly.diff} simulates column phases that JAX traces (a \code{DiffModel} of phase
functions, each with a \code{differentiable} flag) and differentiates their outputs with
respect to the scenario's parameters. Normal and uniform draws are pathwise and exact; a Bernoulli choice
goes through the straight-through estimator \citep{bengio2013estimating}, Gumbel-softmax
\citep{jang2017categorical, maddison2017concrete} or the score function
\citep{williams1992simple}, and each warns that the gradient is an approximation.
\code{check\_gradients} compares a gradient with central finite differences in double
precision, and \code{GradientCalibration} runs Adam \citep{kingma2015adam} over
\code{modelly.calibrate}'s priors and targets, deterministic for a seed.

\paragraph{Measured} (\code{tests/diff}, \code{docs/differentiable.md}), against runs of the
agent-based zoo models at known parameters: the differentiable twin of Bass diffusion (20,000
households, 100 steps, straight-through) recovers an innovation rate of 0.0305 and an
imitation rate of 0.379 for a truth of 0.03 and 0.38, within 2\% (the test allows 5\%); the
Bouchaud--Mézard twin (1,000 agents, 300 steps, two replicates, pathwise) recovers a
volatility of 1.448 and an exchange rate of 0.937 for 1.5 and 1.0, within 7\% (the test allows
10\%), and 1.54 and 1.09 with four replicates, so the error is the observed Gini series'
noise, not the gradient's. A pathwise gradient agrees with finite differences to a relative error below $10^{-6}$ on a
one-parameter test model and below $10^{-4}$ on the Bouchaud--Mézard twin, and the
Gumbel-softmax Bass model within $10^{-3}$; the straight-through
estimator fails the same check, by design, since it is biased. The tests skip without JAX;
on the GPU environment all 105 GPU and differentiable tests pass on CuPy and CUDA JAX.

\subsection{Compiled event heap, column gather and memory-mapped columns}
\label{app:perf2}

A fifth C module, \code{fastevents}, holds the event list's heap: the same sift steps as
\code{heapq}, on keys (time, priority, insertion order) read once at push and compared as C
numbers, with a line-for-line Python twin that gives the same pops and the same heap array on
random schedules, and checkpoints that load in either (\code{tests/native/test\_event\_heap*.py}).
The data collector's attribute reporters, \code{AgentSet.get} and \code{agg} read a
\code{Column} attribute from its table in C instead of calling the Python getter per agent,
with identical rows (\code{test\_column\_gather.py}). \code{MemoryMapped(dir)} keeps a
table's columns in \code{numpy.memmap} files, grown in place, with checkpoints that carry the
rows. A cohort table of 100 million rows (19 bytes a row) ran three life-cycle steps inside a 16~GB
\code{systemd} memory limit with a peak of 3.61~GB memory-mapped (4.27~GB in memory); a step
took 10.4 to 12.0~s memory-mapped and 9.8 to 10.5~s in memory (one run of three steps each);
and under a 3~GB limit the memory-mapped table finished, its steps taking 10.5 to 13.7~s,
while the in-memory one was killed (\code{benchmarks/large\_columns.py}); 100 million \emph{agent objects} cannot fit
in 16~GB, since the objects alone need over 22~GB.

\subsection{Live data and data assimilation}
\label{app:assimilation}

\code{modelly.stream} hands observations to a running model as they arrive: a replayed table,
a CSV file that grows between polls, an HTTP endpoint, and MQTT or Kafka topics through
optional extras (tested against fake clients only; no broker was available). An
\code{ObservationOperator} reads the observed quantities from a model and scores an
observation under independent Gaussian errors. An \code{Ensemble} advances $n$ replicates of
one scenario in lock step in one process, with spawned seeds and starting values drawn from a
prior; a \code{ParticleFilter} is the bootstrap filter of \citet{gordon1993novel}: it weights
the replicates by the likelihood, tracks the effective sample size, and resamples
systematically \citep{douc2005comparison} below a threshold, duplicating replicates through
checkpoint copies reseeded from derived keys; an \code{EnsembleKalmanFilter} is the
stochastic filter with perturbed observations \citep{evensen1994sequential, evensen2003ensemble}
on named model attributes, with optional bounds. Both give the ensemble's mean and quantiles,
now and ahead, as a long-format table.

\paragraph{Measured} (the tutorial \code{docs/assimilation.md}, which a test in
\code{tests/stream/} runs as a module): a town of 1,000 people with a true
contact rate of 0.45 produces 50 days of noisy case counts; 40 replicates with contact rates
drawn from a uniform prior track it. The root-mean-square error of the filtered nowcast of the
unobserved number infected is 5.07 people, against 69.2 for the unfiltered ensemble (ratio
0.073; the test requires below 0.2), with 24 resamplings; the ensemble Kalman filter
recovers a contact rate of 0.451 (spread 0.030 across replicates) for the true 0.45 (the test
allows $\pm 0.05$). On a linear-Gaussian case with 2,000 replicates the Kalman update
reproduces the exact posterior. The tutorial uses one synthetic dataset (truth seed 1,
ensemble seeds 2 and 3); a coverage study over datasets was not done. Two runs from one seed
give identical weights, histories and
replicate states, in one process and in a fresh interpreter, with the compiled and the
pure-Python core. Stepping stays in one process (shipping every replicate to a worker for each
observation would cost more than it saves), and in a model with per-agent streams a
duplicate's agents keep the streams of the agents they were copied from, which the
documentation states. The numbers of this paragraph are recorded in
\code{paper/data/assimilation.json} by \code{paper/scripts/assimilation.py}.

\subsection{Fields, differential equations and a coupling schedule}
\label{app:couple}

\code{modelly.couple} couples agents to continuous processes. A \code{Field} owns or adopts a
float property layer of an orthogonal grid and updates it in place with diffusion, first-order
upwind advection and exact decay, in explicit sub-steps that refuse a time step above the
stability limit (pointing to the implicit scheme, backward Euler with a sparse LU solve), with
closed, absorbing or periodic boundaries and a budget of what was added, flowed out and
decayed. \code{ODE} and \code{StockFlow} components integrate with \code{solve\_ivp}. A
\code{Coupling} orders components phase by phase, writers before readers, and \code{build()}
rejects a read nobody writes (naming the closest quantity), two writers of one quantity, a
cycle inside a phase, and a lagged read without an initial value.

\paragraph{Measured} (\code{tests/couple}): closed and periodic boundaries conserve mass to a
relative $10^{-12}$ in both schemes; a spreading pulse's variance grows by exactly $2Dt$ per
axis; a Gaussian matches the analytic heat kernel, with an error near $10^{-6}$ at $D\,\Delta
t/h^2 = 1/6$, where the explicit scheme is fourth-order; upwind drift moves the centre of mass
by exactly $vt$; decay follows $e^{-kt}$; with absorbing boundaries what remains plus what
flowed out equals the initial mass; logistic growth matches its closed form for four
integrators; and a model with an implicit field and an ODE saves and loads bit for bit. Two
zoo models use it. In \code{nutrient\_foragers} (a nutrient entering at a river bank and
diffusing inland, eaten by foragers) the mean population is within 3\% of the carrying
capacity from the energy balance, within 1\% once the measured decay and starvation losses
are counted, the nutrient budget closes to round-off, and forager density falls with the
distance from the river. In \code{hospital\_epidemic} (SIR agents whose admissions drive a
hospital stock-flow ODE that damps contacts) the attack rate matches the final-size relation
within five standard deviations without the contact response, occupancy and beds match the
exact matrix-exponential solution within $10^{-6}$, admissions equal patients in hospital plus
discharged to round-off, and a stronger response lowers both the peak and the attack rate.
Fields work on orthogonal grids with uniform spacing and coefficients (a hexagonal grid is
refused), the advection scheme smears sharp fronts, and since decay uses \code{exp} and
\code{solve\_ivp} chooses its steps from an error norm, these two outputs are compared across
platforms with a tolerance (\cref{sec:reproducibility}).

\subsection{Emulators, active learning and history matching}
\label{app:emulate}

\code{modelly.emulate} trains a Gaussian-process emulator (the process of
\code{modelly.calibrate}, with restarts, covariances and conditioning) or a \code{numpy}-only
ensemble of small networks on an experiment's runs; \code{ActiveLearning} picks the next runs
by integrated variance reduction (the default) or largest variance \citep{mackay1992information},
in batches conditioned on each pick; \code{HistoryMatch} \citep{vernon2010galaxy,
andrianakis2015bayesian} rules out the parameter space whose implausibility exceeds 3, in
waves, refitting the emulators inside the space that remains.

\paragraph{Measured} (\code{tests/emulate}, \code{docs/emulators.md}). History matching of the
age-structured SEIR zoo model over three parameters, from three observed outputs (peak share
infectious, its timing, final size) at the zoo defaults: four waves of 20 runs leave 0.1705,
0.0444, 0.0307 and 0.0246 of the box, the truth inside after every wave (its largest
implausibility 2.2), and \code{MODELLY\_PURE\_PYTHON=1} rules out exactly the same points.
Active learning against a Latin hypercube (\code{benchmarks/emulate/active\_vs\_lhs.py}:
normalised root-mean-square error on 500 test points, the hypercube's median of three seeds,
the factor being the hypercube runs needed for the active design's error divided by the
active runs): on the Bass diffusion zoo model the factor is 1.64, 2.64, 1.33 and 2.67 at 10
to 40 runs and above 2.4 at 50 and above 2.0 at 60 (no hypercube of up to 120 runs reached the active
design's error); on Branin \citep{dixon1978towards} 1.13 to 1.43 below
80 runs; on the Ishigami function 0.79 to 1.41. Active learning halves the runs a
hypercube needs (a factor of 2) on the zoo model at most budgets, and not on the two
analytic functions, where a well-spread design is already close to the best a
Gaussian-process variance, which does not depend on the outputs, can do; the largest-variance
criterion (the alternative to integrated variance reduction) is worse than a hypercube on
Ishigami at four of six budgets (0.20 to 0.40) and better at 20 runs (1.82). The tests assert what was
measured (40 active Branin runs match a hypercube of 50; 20 active Bass runs match 40).

\subsection{Synthetic populations and models on real data}
\label{app:population}

\code{modelly.population} builds a starting population from aggregate tables: iterative
proportional fitting \citep{deming1940least} with truncate-replicate-sample integerisation
\citep{lovelace2013truncate}, simulated annealing over whole records of a sample, households
from a size table with an adult head, households placed in regions, schools and workplaces
assigned to the nearest place with room within a distance limit, and a contact network with
household, school and work layers; the age-structured SEIR zoo model can start from it. No
personal data is bundled: the marginals of the tests are made up. Three small samples of
Burlington, Vermont, ship with their licences in a \code{NOTICE} file: OpenStreetMap streets
(ODbL 1.0, © OpenStreetMap contributors; 255 junctions, 607 directed segments), a $120
\times 120$ window of ESA WorldCover 2021 \citep{zanaga2022worldcover} (CC BY 4.0) and 41
census tracts (public domain), with a script that downloads the full sources, checks their
SHA-256 and rebuilt the samples byte for byte.

\paragraph{Measured} (\code{tests/population}, \code{tests/zoo}): about 20,000 synthetic people
over the 41 tracts reproduce their totals and age counts exactly in every tract, their sex
counts within fewer people than there are age levels (under 1\% in total), household-size
shares within 0.02, every household with an adult, no school or workplace over its 600
seats or beyond its distance, and the SEIR model started from them follows the final-size
relation within 0.05. Three zoo models run on the real samples: \code{road\_traffic}, whose
shortest-path flows equal NetworkX's edge betweenness on all 607 segments and stay within
five binomial standard deviations over three seeds; \code{habitat\_selection}, whose
occupancy per land-cover class reaches its Metropolis stationary law within five standard
deviations and never enters water; and \code{schelling\_regions}, whose dissimilarity at a
threshold of 0.45 ends more than five standard deviations above random allocation and rises
with the threshold.

\subsection{Modelly in the browser}
\label{app:browser}

\cmd{modelly export-web model.py:Model} writes a static site that runs the model with Pyodide
\citep{pyodide} (pinned version 314.0.7, loaded from its content delivery network): the parameter panel of Studio,
run controls, a space view, a reporter plot, the model in a Web Worker so the page answers
during long runs, Modelly as a trimmed pure-Python wheel of about 0.5~MB, and only the
packages the model was observed to import. The C modules also build for WebAssembly with
Emscripten through \code{pyodide build}, and the export can ship that wheel.

\paragraph{Measured} (\code{tests/web/test\_browser.py}, \code{docs/browser.md}). The 14
bundled examples (the three split-run reference models were not exported) and
the 17 zoo models, 31 models, each run from its export in headless Chromium with the pure
wheel and with the compiled one, 62 of 62 runs without error. Because Pyodide's GDAL reads no Deflate-compressed GeoTIFF and
its geopandas has neither pyogrio nor fiona, the bundled land-cover window is stored
uncompressed and the tracts are read through \code{GeoDataFrame.from\_features}. A cold
first visit transfers 14.3~MB for the age-structured SEIR model
(Pyodide, its standard library, \code{numpy} and \code{pandas} are 13.6~MB of it) and is
ready at step 0 after 4.4~s unthrottled, 7.0~s on a 50~Mbit/s line with 40~ms round trips, and
20.0~s on Chrome's ``Fast 4G'' profile (about 1~MB/s, 165~ms), where the transfer alone takes
14~s; the forest fire, a model that needs \code{scipy} (as the grid, continuous and network
spaces do), transfers 28.3~MB and takes 5.4~s unthrottled, 10.1~s on the 50~Mbit/s line and
35.0~s on Fast 4G (medians of three runs; \code{tools/web/measure\_load.py}). A first load
under ten seconds on a good home connection is reached without \code{scipy}, missed narrowly
with it, and not reached on a mobile profile.
The documentation gallery embeds three live models.
